\chapter{总结与展望}

本文围绕江南大学博世·飞驰人生队的飞跃雷区空地协同系统，梳理了机械与硬件平台、双核实时软件、
姿态与位置控制、视觉定位、多传感器融合、空地通信以及调试验证方法。当前仓库已经具备较完整的无人机端实现与工具链，
包括CYT4BB7双核分工、ICM42688姿态链路、双ToF与光流速度观测、双摄像头视觉处理、
多级闭环控制、DShot600、电调遥测、VOFA诊断和RotorPy仿真。

现阶段仍需围绕上场定型实物补全机械与电气参数、完整原理图、车模端实现说明及全链路量化实验。
后续工作重点是以统一数据集和重复试验评估定位精度、控制误差、通信时延、任务成功率和完成时间，
并将全部结论关联到明确的代码版本、标定文件和原始数据。

\todoitem{定稿时重写本章：概括已经取得的具体指标，客观列出仍存在的问题及改进方向；不要重复摘要，也不要引入正文未验证的新结论。}
